\section{Crossing near a rationally chain connected set}\label{app:crossing}

The rational boxes now give a crossing map on a neighborhood of an
entire projective set. The passage has two tasks: provide continuation
along enough paths in the set, and remove path dependence while retaining
ambient holomorphic germs. In this appendix \(M\) is smooth and projective, and
\(T_M=E_1\oplus E_2\) is an integrable splitting. A set is
\emph{joined by rational chains} if every two of its points can be
joined inside it by a finite chain of images of \(\Pp^1\).

\begin{theorem}[Neighborhood crossing]\label{thm:crossing}
Let \(P\subset M\) be a nonempty connected projective algebraic set.
Suppose that \(P\) is simply connected and joined by rational chains.
There are a connected open neighborhood \(U\) of \(P\) and a holomorphic
map \(c:U\times U\to M\) that agrees near the diagonal with ordinary
crossing in product charts. Throughout \(U\times U\), it satisfies
\begin{equation}\label{cross:identities}
\begin{gathered}
 c(x,x)=x,\qquad
 \dd_1c(T_M)\subset E_1,\quad \dd_1c(E_2)=0,\\
 \dd_2c(T_M)\subset E_2,\quad \dd_2c(E_1)=0.
\end{gathered}
\end{equation}
The image subbundles in these formulas are evaluated at \(c(x,y)\).
\end{theorem}

\subsection{Continuation of ambient germs}

Put \(Q=P\times P\), fix \(p\in P\), and write \(q_*=(p,p)\).
All germs below are germs of holomorphic maps from \(M\times M\) to \(M\),
even when their basepoints lie in \(Q\).

\begin{definition}\label{cross:attained}
A germ over \((a,b)\in Q\) is \emph{attained} if it is the crossing germ
of Lemma~\ref{rat:germs} associated to a rational-tree box \(R_u\) of
Proposition~\ref{rat:boxes}, at two distinct smooth marked points \(s,t\) of
its parameter tree, with \(u(s)=a\) and \(u(t)=b\).
\end{definition}

The ordinary crossing at \(q_*\) is attained using a constant sphere.
We do not assume that different realizing trees give the same germ.

\begin{lemma}\label{cross:tree-path}
Let \(k:D\to Q\) be a holomorphic map from a rational tree with distinct
smooth markings \(m_1,m_2\). Every attained germ at \(k(m_1)\) continues
along the image of any path from \(m_1\) to \(m_2\). Its terminal germ
is attained and is independent of the path in \(D\).
\end{lemma}
\begin{proof}
Choose a realizing tree \(T\) with markings \(s,t\) for the initial
germ. Attach two separate copies of \(D\) to \(T\), identifying their
first markings with \(s\) and \(t\), respectively. Map the copies to
\(M\) by the two components of \(k\). The resulting curve is again a
rational tree, and the maps agree at the attachment points. Its box
restricts to the original box, including its ambient germs, by the
uniqueness and branch compatibility in Proposition~\ref{rat:boxes}
and Lemma~\ref{rat:germs}.

Move the two parameter points along the copies of a path in \(D\).
Lemma~\ref{rat:germs} gives exactly the asserted continuation along
its image under \(k\). Its terminal markings are smooth and distinct,
so the terminal germ is attained. All paths between the markings are
homotopic in the simply connected tree; the corresponding leafwise
holonomy is unchanged.
\end{proof}

We will use the following elementary stability fact. It applies to
holomorphic germs on any complex manifold, with values in another complex
manifold.

\begin{lemma}[Stability of continuation]\label{cross:stability}
Suppose that a germ continues along a continuous path \(\gamma:[0,1]\to N\)
in a complex manifold \(N\). There are a representative of its terminal
germ and a uniform neighborhood of \(\gamma\) in the space of paths
starting at \(\gamma(0)\) such that continuation along every path in that
neighborhood exists and ends in that representative.
\end{lemma}
\begin{proof}
Choose a finite subdivision and holomorphic representatives on open
neighborhoods of the compact subpaths. Consecutive representatives agree
on an open neighborhood of each subdivision point. After shrinking the
subpath neighborhoods and these overlaps, every sufficiently uniformly
close path stays in the same domains and crosses through the same
equality overlaps. The same finite chain of representatives works.
\end{proof}

Continuation along a fixed path from a fixed germ is unique. Indeed,
two continuation chains that agree initially agree locally at the first
possible separation point, by the identity theorem in overlapping
ambient neighborhoods. This also proves that continuation along a path
followed by its reverse returns the initial germ.

\subsection{A compact family of rational routes}

Embed \(Q\) in a projective space \(\Pp^N\). We use the following standard
properties of genus-zero stable maps: fixed-degree spaces with three
markings have projective coarse moduli spaces; evaluations are algebraic;
maps factoring through a fixed projective subscheme form a closed locus;
and locally the projective-space moduli problem has finite-quotient charts
with universal families of nodal curves. These statements follow from
Fulton--Pandharipande \cite[Theorem~1, Proposition~3, and Sections~1.3, 4.1, 4.4, and 5.1]{FP97}.
For the local families we always use the ambient target \(\Pp^N\), which
is convex, and restrict to the locus of maps into \(Q\).

\begin{lemma}\label{cross:parameter-space}
There is a compact semialgebraic space \(S_*\) of three-pointed
rational-tree maps into \(Q\) whose first evaluation is \(q_*\) and
whose second evaluation
\[
 e:S_*\longrightarrow Q
\]
is surjective.
\end{lemma}
\begin{proof}
For each degree \(d\ge0\), take the closed locus \(S_d\) in
\(\overline M_{0,3}(\Pp^N,d)\) of maps into \(Q\) with first evaluation
\(q_*\). Its second evaluation image \(F_d\) is closed algebraic.
Every point of \(Q\) lies in some \(F_d\): join it to \(q_*\) by moving
first in one factor and then in the other along rational chains.
Represent the chain by a rational tree. Insert contracted components if
needed to put the endpoints at distinct smooth markings, add a third
marking, and stabilize. This does not change the endpoint evaluations.

Write \(Q=Q_1\cup\cdots\cup Q_r\) for its irreducible components.
For each \(i\),
\[
 Q_i=\bigcup_{d\ge0}(Q_i\cap F_d).
\]
If none of the \(F_d\) contained \(Q_i\), this would express the compact
complex variety \(Q_i\) as a countable union of closed nowhere dense
sets, contradicting the Baire category theorem. Choose \(d_i\) with
\(Q_i\subset F_{d_i}\). The disjoint union of the finitely many spaces
\(S_{d_i}\), omitting repetitions, is the required \(S_*\).
\end{proof}

\begin{lemma}[Routes in a degenerating tree]\label{cross:nearby-routes}
Fix a stable-map parameter in \(S_*\) and a path between its first two
markings, chosen along smooth component arcs and through the intervening
nodes. For every sufficiently nearby parameter one can choose a path
between its corresponding markings whose image in \(Q\), after
parametrization on \([0,1]\), is uniformly as close as desired to the
fixed image path.
\end{lemma}
\begin{proof}
Choose a representative of the limiting stable map in a local
finite-quotient chart with a family. A finite quotient map is open, so
nearby coarse parameters admit representatives close to this one;
no continuous choice is required.

Follow the finite component route between the markings. Away from
small segments around its nodes, finitely many relative smooth charts
move the chosen arcs to nearby fibers, with the prescribed marked ends.
At a node the family is locally \(xy=a\), where \(a\) is a function
of the parameter. The incoming and outgoing nearby pieces can be joined
within an arbitrarily small plumbing neighborhood when \(a\) is small;
the same assertion holds for \(a=0\). The family map is continuous,
so all images in that neighborhood lie close to the limiting nodal
value. Parametrize these joining pieces on the corresponding small
subintervals. The resulting paths have uniformly close images.
Contracted components cause no difficulty, since their limiting images
are constant. Restricting the parameter locus to maps into \(Q\)
keeps all these image paths inside \(Q\).
\end{proof}

\begin{lemma}\label{cross:semialgebraic-lift}
Every continuous semialgebraic path \(\sigma:[0,1]\to Q\) has a
finite subdivision such that on each closed subinterval it admits a
continuous lift to \(S_*\). Lifts on consecutive subintervals need not
agree at their common endpoint.
\end{lemma}
\begin{proof}
Semialgebraic choice applied to
\[
 \{(t,s)\in[0,1]\times S_*:e(s)=\sigma(t)\}\longrightarrow[0,1]
\]
gives a semialgebraic section. It is continuous on finitely many open
subintervals. Compactness of \(S_*\) and one-variable semialgebraic
monotonicity give one-sided limits at every endpoint. Continuity of
\(e\) makes those limits lift the correct endpoint values. Extend each
interval lift by its one-sided limits. We use the standard semialgebraic
choice and one-variable finiteness results in
\cite[Theorems~2.1, 3.1, and 3.4]{Coste99}; see also \cite{BCR98}.
\end{proof}

\subsection{Monodromy and an actual neighborhood}

\begin{lemma}\label{cross:all-paths}
Every attained germ continues along every continuous semialgebraic path
in \(Q\) starting at its basepoint, and the terminal germ is attained.
\end{lemma}
\begin{proof}
Use Lemma~\ref{cross:semialgebraic-lift} and consider one lifted
interval \([a,b]\), with prescribed attained germ \(g_a\) at
\(\sigma(a)\). Continue \(g_a\) backwards along a marked route in the
tree at \(a\), obtaining an attained germ \(h\) at \(q_*\) by
Lemma~\ref{cross:tree-path}. For each \(t\in[a,b]\), continue this
\emph{fixed} germ \(h\) forwards through the tree over \(t\), and denote
the endpoint germ by \(g_t\). The choice of route does not affect it.
At \(t=a\), reverse-path uniqueness gives the prescribed \(g_a\).

For \(t\) near any \(t_0\), Lemma~\ref{cross:nearby-routes} permits
uniformly close image routes from the common point \(q_*\).
Lemma~\ref{cross:stability} then supplies one ambient representative
giving all \(g_t\) near \(t_0\). Thus these germs are analytic
continuation along \(\sigma\) on the whole interval. Each is attained
by Lemma~\ref{cross:tree-path}. Repeat from the terminal germ on the
next interval. The intermediate germ \(h\) at \(q_*\) is not assumed
to be the original ordinary crossing; only its being fixed on the
current interval is used.
\end{proof}

\begin{proof}[Proof of Theorem~\ref{thm:crossing}]
The compact family of routes has supplied continuation along every
semialgebraic path, including through singularities and changes of
component. It remains to identify the terminal germs and glue them.
Triangulate \(Q\) semialgebraically
\cite[Theorem~4.4]{Coste99}. Since \(P\) is simply connected,
so is \(Q\). Start with the ordinary crossing germ at \(q_*\), and
continue it along piecewise linear paths in the triangulation. These
paths are semialgebraic in \(Q\), so Lemma~\ref{cross:all-paths}
applies. Any two such paths with the same endpoints are homotopic
relative endpoints through paths with semialgebraic slices: apply
relative piecewise linear approximation to a nullhomotopy in the
finite triangulation. For a fixed slice, compact-path stability
(Lemma~\ref{cross:stability}) and uniform continuity of the homotopy
show that nearby slices have the same terminal germ. The terminal
germ is consequently constant throughout the homotopy. We obtain
well-defined ambient germs at every point of \(Q\).

These germs are locally coherent along \(Q\). Indeed, choose a small
relative path-connected neighborhood of an endpoint inside the domain
of its representative. Appending small paths that are piecewise linear in the fixed
semialgebraic triangulation, after subdivision, and lie in that
neighborhood shows that the same representative gives the assigned
germs there. Such neighborhoods are available from the triangulation.

For completeness, this coherent family has a representative on an
actual neighborhood of \(Q\). Choose finitely many representatives
\(c_i:V_i\to M\) giving the family on \(Q\cap V_i\), and smaller open
sets \(W_i\) with compact closures in \(V_i\), still covering \(Q\).
On \(V_i\cap V_j\), the set where the germs of \(c_i,c_j\) agree is
a union of connected components, by the identity theorem. Its
complement is therefore relatively closed. The compact sets
\[
 B_{ij}=\{x\in\overline W_i\cap\overline W_j:
                 (c_i)_x\ne(c_j)_x\}
\]
are disjoint from \(Q\). Removing their finite union and restricting
to \(\bigcup_iW_i\) gives a neighborhood of \(Q\) on which the
representatives glue holomorphically. Notice that equality here is
equality of \emph{ambient germs}, not merely of values on \(Q\).

Compactness of \(P\) gives an open neighborhood \(U\) of \(P\) whose
square lies in this domain. Since \(P\) is connected, take the connected
component of \(U\) containing it. The glued map restricts to
\(c:U\times U\to M\). Its differential identities in
\eqref{cross:identities} hold near \((p,p)\) and hence throughout
connected \(U\times U\) by holomorphicity. The identity on the diagonal
extends throughout connected \(U\). In a product chart about any
\(x\in U\), these differential identities make the first output
coordinate depend only on the first coordinate of the first input,
and the second output only on the second coordinate of the second
input. The diagonal identity identifies both functions with the
corresponding coordinates. Hence \(c\) is ordinary crossing near every
diagonal point, as required.
\end{proof}
